\begin{subsection}{Lattice Topologies}
    As stated before we loose Lipschitz continuity, if we complete the risk measure lattice,
    since expressions such as \( | \rho_{\top}(X) - \rho_{\top}(Y) | \) are undefined.
    Generally risk measures are, as far as the author is informed, not typically studied under the aspect
    of their own topology. 
    \newline
    The extended real numbers \( \extreals \) are second countable, compact, metrizable and \( \mathbb{R} \) is dense in \( \extreals \) 
    (Example 2.75 in \cite{infinitedimanalysis}).
    One way (there are more) to grant them a metric is by noting that there is a homeomorphism 
    \[\psi : \extreals \to [0, 1], x \mapsto \frac{x}{1 + |x|} \]
    and then to define \( d_{\extreals}(x, y) = d_{\mathbb{R}}(\psi(x), \psi(y)) = |\psi(x) - \psi(y)| \).
    Similarly we can proceed to endow \( \extrm \) with a metric.
    \begin{theorem} \label{unifconvrm}
        One can define a metric on \( \extrm \) via
        \[
        \dextrm{\rho_1}{\rho_2} = \sup_{X \in \mathcal{X}} d_{\extreals}(\rho_1(X), \rho_2(X)),
        \]
        where \( d_{\extreals}(\cdot, \cdot) \) is the metric of the extended real numbers.
        \( \extrm \) is with respect to the induced topology a complete space 
        in which \( \mathcal{R} \) is dense.
    \end{theorem}
    \begin{proof}
    Clearly \( \dextrm{\rho}{\rho} = 0 \) and \( \dextrm{\rho_1}{\rho_2} = \dextrm{\rho_2}{\rho_1} \), since \( d_{\extreals}(\cdot, \cdot) \) is a metric.
    Further if \( \dextrm{\rho_1}{\rho_2} = 0 \), then \( \rho_1(X) = \rho_2(X) \) for all \( X \in \mathcal{X} \) (as \( \psi \) is injective),
        and \( \rho_1 = \rho_2 \). Moreover,
        \begin{align*}
            \sup_{X \in \mathcal{X}} d_{\extreals}(\rho_1(X), \rho_3(X)) 
            &\leq \sup_{X \in \mathcal{X}} d_{\extreals}(\rho_1(X), \rho_2(X)) + d_{\extreals}(\rho_2(X), \rho_3(X)) \\
            &\leq \sup_{X \in \mathcal{X}} d_{\extreals}(\rho_1(X), \rho_2(X)) + \sup_{X \in \mathcal{X}} d_{\extreals}(\rho_2(X), \rho_3(X)) \\
            &= \dextrm{\rho_1}{\rho_2} + \dextrm{\rho_2}{\rho_3},
        \end{align*}
            where the triangle inequality of \( d_{\extreals}(\cdot, \cdot) \) is applied pointwise for the first inequality.
            Thus \( \dextrm{\cdot}{\cdot} \) is a metric.
            \newline
            Take a convergent sequence \( \{ \rho_n \}_{n \in \mathbb{N}} \subset \extrm \), 
            and any \( X, Y \in \mathcal{X} \) s.t. \( X \leq Y \), for each \( n \) it holds
            \begin{align*}
                \rho_n(Y) &\leq \rho_n(X) \\
                \psi(\rho_n(Y)) &\leq \psi(\rho_n(X)) \\
                \lim_{n \to \infty} \psi(\rho_n(Y)) &\leq \lim_{n \to \infty} \psi(\rho_n(X)) \\
                \psi(\lim_{n \to \infty} \rho_n(Y)) &\leq \psi(\lim_{n \to \infty} \rho_n(X)) \\
                \psi(\rho(Y)) &\leq \psi(\rho(X)) \\
                \rho(Y) &\leq \rho(X)
            \end{align*}
            due to \( \psi \) being monotone, continuous, invertible and non-strict inequalities being preserved by taking limits.
            As \( \rho(X + m) = \lim_{n \to \infty} \rho_n(X + m) = \lim_{n \to \infty} \rho_n(X) - m = \rho(X) - m \),
            \( \rho \in \extrm \), whence \( \extrm \) is complete.
            Consider the sequence \( \{ \rho_n \}_{n \in \mathbb{N}},\ \rho_n = \rho + n \) for some \( \rho \in \mathcal{R} \),
            clearly \( \lim_{n \to \infty} \rho_n = \rho_\top \) and similarly we can construct a sequence for \( \rho_\bot \)
            hence the closure of \( \mathcal{R} \) is \( \extrm \).
    \end{proof}
    As in the case of the continuous functions, the metric from Theorem \ref{unifconvrm} can be interpreted as being induced by uniform convergence
    of risk measures. As opposed to robustness in regard to the probability measure, risk measures which are close in this space, are in so far
    robust that it shouldn't matter much, which risk measure we use.
    The application of this analysis is that we can ensure that our lattice structure is compatible with a desirable topology 
    on \( \extrm \).
    \begin{lemma} \label{maxminarecont}
        The \( \max, \min : \extreals \times \extreals \to \extreals \) are continuous.
    \end{lemma}
    \begin{proof}
        We show the result for \( \max \), the proof for \( \min \) is similar.
        A subbasis for the topology of \( \extreals \) is given by all intervals of the form \( (x)^\upuparr \)
        and \( (x)^\dodoarr \) (see Example 2.2.6 in \cite{infinitedimanalysis}).
        We have \[ 
            (\max)^{-1} ((x)^\upuparr) = ((x)^\upuparr \times (x)^\downarrow) \cup ((x)^\downarrow \times (x)^\upuparr),
        \]
        and as each part of the union contains the open set \( (x)^\upuparr \) so is its product with any set
        and the union of two open sets is open as well.
    \end{proof}
    \begin{corollary}
        The lattice \( \extrm \) is a topological lattice.
    \end{corollary}
    \begin{proof}
        Since \( \extreals \) is compact, any continuous function on $\extreals \times \extreals$ is 
        uniformly continuous, in particular the \( \max, \min \) by Lemma \ref{maxminarecont}.
        Therefore, for every \(\epsilon > 0\), there exists a \(\delta > 0\) such that for any \(a, b, a', b' \in \extreals\):
        \[
        d_{\extreals}(a, b) < \delta \quad \text{and} \quad d_{\extreals}(a', b') < \delta 
        \]
        so by continuity \( d_{\extreals}(\max(a, a'), \max(b, b')) < \epsilon \).
        \newline

        Assume \(d_{\extrm}(\rho_{0}, \rho_{1}) < \delta\) and \(d_{\extrm}(\rho_0', \rho_1') < \delta\). 
        By the definition of \( \dextrm{\cdot}{\cdot} \), this implies that for all \(X \in \mathcal{X}\):
        \[
            d_{\extreals}(\rho_{0}(X), \rho_1(X)) < \delta \quad \text{and} \quad d_{\extreals}(\rho_0'(X), \rho_1'(X)) < \delta
        \]

        Using the uniform continuity from \( \extreals \), it follows that for all $X \in \mathcal{X}$:
        \[
            d_{\extreals}\big(\max(\rho_0(X), \rho_0'(X)), \max(\rho_1(X), \rho_1'(X))\big) < \epsilon.
        \]
        Taking the supremum over all \( X \in \mathcal{X}\):
        \[
            d_{\extrm}(\rho_0 \vee \rho_0', \rho_1 \vee \rho_1') \le \epsilon,
        \]
        and the result holds vis-à-vis for \( \wedge \).
    \end{proof}

   
    Now we can also give some topological classification to some algebraic results from before.
    For example Lemma \ref{necessaryhilbert} now also requires that any ideal \( P \)  
    in \( \extrm \) for which it holds \( \acceptidealb{\riskidealb{P}} = P \), is a closed convex set in \( \extrm \),
    meaning it is the intersection of the kernel of the linear functions \( f_X : \extrm \to \extreals, \rho \mapsto \rho(X) \)
    for \( X \in S \).
    %This could potentially be an analogue to the concept of convex ideals from real algebraic geometry (Definition \ref{defconvexideal}).
\end{subsection}
